routing data. Sections 93--96 bound or rank the canonical $A$-streak and the
subsequent B-selecting transfers. Sections 94--95 rank factor forks, and
Sections 97--135 transport the finite tokens through every high-return and
long-tail recurrence. Every operation that changes a retained source or token
is supported by a strict comparison; pure routing preserves the retained
projection and is well-founded inside the fixed fibre.

\begin{lemma}[Well-founded fibre principle]\label{lem:fibre}
Let a transition relation have an ordinal-valued projection $\pi$. If
$C\to C'$ implies $\pi(C')\le\pi(C)$ and the relation restricted to every
fibre $\pi^{-1}(\{\alpha\})$ is well-founded, then the full relation is
well-founded.
\end{lemma}

\section{The external audit and the repaired entry}\label{sec:audit}

An external audit of an earlier version found two genuine defects.

First, the proof used
\[
 \text{smaller canonical coefficient}
 \Longrightarrow
 \text{smaller coefficient source},
\]
which is false when powers of three are present. The explicit counterexample
is
\[
 s=1,\qquad J(1)=5,\qquad a=3J(1)=15,\qquad e=1.
\]
The common boundary child is
\[
 R(3a-e)=R(44)=22.
\]
Its canonical coefficient is $11<15$, but $11=J(3)$, so the source has
increased from $1$ to $3$.

Second, the reverse-factor lemma was used at exponent one even though it is
valid only from exponent two. Hence the family
\[
 Q_1^e(3^kJ(s)),\qquad k>0,
\]
needed a separate entry proof. Section~\ref{sec:audit} gives that proof
without reinstating either invalid shortcut.

\subsection{One-shot initialization}

The key observation is that there are two logically different operations:
initializing the typed normalizer after the raw entry has been normalized,
and resetting a normalizer that is already running. Only the second requires
a preceding strict source/token decrease.

Refine the retained lexicographic rank by an entry bit
\[
 \eta\in\{1,0\},\qquad 1>0,
\]
placed after the outer source and outer token multiset:
\[
 \boxed{
 (\text{outer source},\ \PhiM(M),\ \eta,\
   \text{inner source/token rank},\ \text{finite control rank}).}
\]
Here $\eta=1$ means that the current outer fibre has not yet initialized its
typed inner normalizer, and $\eta=0$ means that Proposition
\ref{prop:typed-normalizer} is running.

The transition
\[
 \boxed{\eta:1\longrightarrow0}
\]
is strict independently of the numerical value assigned to the newly
initialized inner source. Thus the first returned source may be larger than
the original source. Within an unchanged outer source/token fibre, the bit can
never return from $0$ to $1$. A later raw reinitialization is allowed only
after an earlier outer source or outer proof-token component has strictly
decreased, exactly as permitted by Lemma~\ref{lem:fibre}. Therefore the entry
bit cannot hide a repeated source-growth cycle.

\subsection{The factor-free exponent-two bootstrap}

Let $x>0$, $a=J(x)$, and put
\[
 P=Q_1^e(a),\qquad U=Q_2^e(a),
\]
\[
 b=B(P)=B(U),\qquad F=A(U)=Q_1^e(3a),\qquad g=1-e.
\]
Suppose $U$ is \DRAW\ at the globally least outer DRAW source $x$. The common
child $b$ has smaller canonical coefficient than $J(x)$, hence its coefficient
source is below $x$; so $b$ is not \DRAW. A child of a \DRAW\ cannot be
\LOSS, therefore
\[
 \boxed{b\text{ is \WIN},\qquad F\text{ is \DRAW}.}
\]
The exact first-factor identity is
\[
 \boxed{9a+1-2e=2^vJ(b)}.
\]
It follows that
\[
 A(F)=Q_v^g(J(b)).
\]
We now split according to whether the input phase at $x$ is B-selecting or
A-selecting.

\subsubsection{B-selecting input}

Assume $e=1-\alpha(x)$. The first source-boundary numerator is divisible by
$4$. The consecutive numerator relation gives
\[
 9J(x)+1-2e
 =3(3J(x)+1-2e)-2(1-2e)
 \equiv2\pmod4,
\]
so
\[
 \boxed{v=1.}
\]
Writing $u=A(x)$ and using the parity that characterizes the B-selecting
phase gives the exact common endpoint
\[
 \boxed{b=3A(x)+1.}
\]
The two children of the \DRAW\ state $F$ are therefore
\[
 T=Q_1^g(J(b)),\qquad C=R(T).
\]
The signed-prefix identity identifies the raw state $C$ with the ordinary
child of $b$ selected by phase $g$:
\[
 \boxed{C\in\{A(b),B(b)\}.}
\]
At least one of $T,C$ is \DRAW.

If $T$ is \DRAW, it is the direct exponent-one typed obligation at source
$b$, and the one-shot bit initializes the normalizer.

If $C$ is \DRAW, the other ordinary child $\ell$ of the \WIN\ state $b$ is
\LOSS. For nonexceptional $b$, the ordinary side diamond exposes
$B(A(b))$ as a \WIN\ child of $\ell$ and as a child of the continuing
\DRAW, hence a strict proof-token/boundary descent. If $b$ is exceptional
and $C=A(b)$, the exceptional side formula gives an adjacent \DRAW\ frame over
the actual \LOSS\ token $\ell$.

The sole remaining exceptional orientation is $C=B(b)$ \DRAW. Since
$b=3A(x)+1$,
\[
 b\le\frac{9x+5}{2}.
\]
For exceptional $b$, the alternating suffix of $A(b)$ has length at least
three, whence
\[
 C=B(b)\le\frac{A(b)}8\le\frac{3b+1}{16}.
\]
If $\rho(C)$ is the coefficient source of $C$, then
\[
 \rho(C)<\frac C6
 \le\frac{3b+1}{96}
 \le\frac{27x+17}{192}<x.
\]
Thus this last row is a genuine strict outer-source transition.

\subsubsection{A-selecting input}

Assume $e=\alpha(x)$. Lemma~\ref{lem:source-transition} gives
\[
 3J(x)+1-2e=2J(A(x)),
\]
and the consecutive numerator identity implies
\[
 \boxed{v\ge2.}
\]
Retain
\[
 9J(x)+1-2e=2^vJ(b).
\]
If $v\ge4$, then
\[
 16J(b)\le9J(x)+1\le27x+19.
\]
Since $J(b)\ge3b+1$,
\[
 48b+16\le27x+19,
 \qquad
 \boxed{b\le\frac{9x+1}{16}<x.}
\]
Both children of $F$ have factor-free source $b$, so this is a strict
outer-source exit.

Only $v=2,3$ remain. If $v=2$, $F$ is a \DRAW\ parent whose children are
\[
 Q_2^g(J(b)),\qquad Q_1^g(J(b)),
\]
which is exactly the second form of the typed obligation $O(b,g)$.

If $v=3$, the children are the valuation-three frame
\[
 Q_3^g(J(b)),\qquad Q_2^g(J(b)).
\]
The constructor automatically supplies the congruence hypothesis needed by
the first-factor hidden-parent lemma. Reducing
\[
 9J(x)+1-2e=8J(b)
\]
modulo $3$ and using $g=1-e$ yields
\[
 \boxed{J(b)\equiv1+g\pmod3.}
\]
Thus that hidden-parent lemma is used only in the typed rows for which its
hypothesis is proved. The ensuing factor-frame normalization produces a
strict source edge, a strict boundary/token edge, or the already typed
obligation entry.

We have proved the following.

\begin{proposition}[Factor-free base-entry attachment]\label{prop:baseattach}
Every factor-free exponent-two \DRAW\ at a globally minimal outer source has
an outcome-compatible finite continuation that either strictly lowers the
retained outer source, strictly lowers a certified finite proof token, or
consumes the one-shot entry bit and enters the typed normalizer.
\end{proposition}

\subsection{Arbitrary factorful exponent one}

Return to
\[
 Q_1^e(3^kJ(s))\text{ \DRAW},\qquad k>0.
\]
The predecessor reduction described after Lemma~\ref{lem:twolevel} has two
outputs. Either it returns to a factor-free \DRAW\ over the same source, in
which case the factor-free analysis terminates at exponent one or at the
exponent-two bootstrap of Proposition~\ref{prop:baseattach}; or it exposes a
factor-free \DRAW\ lift/frame together with an actual finite \WIN/\LOSS\
descendant of the predecessor token. In the latter case Lemma
\ref{lem:multiset} supplies a strict token edge, after which Lemma
\ref{lem:fibre} permits the inner arithmetic task to reset.

Consequently:
\begin{proposition}[Repaired arbitrary exponent-one attachment]\label{prop:attachment}
Every arbitrary factorful exponent-one \DRAW\ has a finite outcome-compatible
normalization that either strictly lowers the retained source, strictly lowers
an already marked finite proof token, or consumes the one-shot entry bit and
enters the typed normalizer. No comparison $t<s$ is required for the temporary
returned inner source.
\end{proposition}

\section{Proof of the main theorem}

\begin{proof}[Proof of Theorem~\ref{thm:main}]
Assume that a \DRAW\ exists and choose the least coefficient source $s$ of a
\DRAW. The source-zero case is handled by the explicit zero-source
constant-tail normalization in the supplement. Assume $s>0$.

The constant-tail/source analysis and the factor scan reduce the continuing
actual \DRAW\ to a marked macro-configuration. By Proposition
\ref{prop:attachment}, every macrostep has one of three effects:
\begin{enumerate}[label=(\arabic*)]
\item the retained outer source strictly decreases;
\item at fixed outer source, the finite proof-token multiset strictly
      decreases under $\PhiM$;
\item with both outer components fixed, the one-shot bit changes from
      $\eta=1$ to $\eta=0$, after which the route lies in the typed fixed-fibre
      normalizer of Proposition~\ref{prop:typed-normalizer}.
\end{enumerate}
The first effect contradicts minimality when it occurs at the initial outer
source and, more generally, is a strict first lexicographic component. The
second is well-founded by Lemma~\ref{lem:multiset}. The third is strict at the
entry bit and cannot reset back to $1$ while the preceding components are
unchanged. Once $\eta=0$, Proposition~\ref{prop:typed-normalizer} gives
well-foundedness of the equal-retained-data fibre. Lemma~\ref{lem:fibre}
therefore makes the complete marked macro relation well-founded.

On the other hand, every \DRAW\ position has no \LOSS\ child and has at least
one \DRAW\ child. Hence from any marked \DRAW\ configuration one may choose an
outcome-compatible continuing \DRAW\ and repeat the finite normalization
macrostep forever. This would produce an infinite path in the well-founded
marked relation, a contradiction.

Therefore the conjugated game has no \DRAW\ positions. Conjugacy with the
original odd-integer game proves that every positive odd starting state has
finite remoteness and optimal play reaches $1$.
\end{proof}

\section{Verification boundary and reproducibility}\label{sec:repro}

The mathematical claim above is a human proof. The computational artifacts
have deliberately narrower roles.

\begin{itemize}
\item The repository's standard-library Python tests check exact arithmetic
      identities, finite outcome certificates, and finite regression ranges.
\item The symbolic global-routing JSON checks the declared typed control/rank
      assembly. It does not encode the new one-shot $\eta$ bridge and therefore
      remains a conditional machine check rather than an end-to-end proof.
\item The Lean project kernel-checks general metatheory and selected definitions,
      but does not kernel-check the complete concrete no-\DRAW\ theorem.
\end{itemize}

The corrected proof package used for this manuscript is the branch
\path{repair/althoefer-audit-gap} of
\url{https://github.com/Grisha-Pochuev/3n-plus-minus-1-game},
with the closure addendum \path{docs/althoefer-audit-closure.md}. At the time
this manuscript was prepared, the corresponding draft pull request was PR~\#1.
The earlier archived preprint is
\url{https://doi.org/10.5281/zenodo.21844684}.
That archived version predates the external audit and should not be used in
place of this corrected manuscript.

The companion file \texttt{verify\_closure.py} supplied with this manuscript
is a standalone arithmetic regression for the new repair identities. It has
no third-party dependencies and explicitly labels itself as supporting
verification, not as a proof.

\section{Audit checklist for the repaired bridge}

An independent reviewer should try to falsify the following points first:
\begin{enumerate}
\item no step identifies a smaller canonical coefficient with a smaller
      coefficient source;
\item the reverse-factor lemma is never used at exponent one;
\item the B-selecting exponent-two base row really has factor valuation one;
\item the exceptional raw B-child estimate gives a source strictly below the
      old source;
\item the A-selecting row with $v\ge4$ really has $b<x$;
\item the $v=3$ constructor really satisfies the additional hidden-parent
